\begin{table*}[htbp]
\centering
\scriptsize
\renewcommand{\arraystretch}{1.35}
\begin{NiceTabularX}{\textwidth}{
>{\centering\arraybackslash}m{0.9cm}
>{\centering\arraybackslash}m{1.8cm}
>{\centering\arraybackslash}m{2.1cm}
X X X X
}[hvlines]
核心维度 & 细分落地 & 核心主体 & 代表性成果 & 技术方案 & 核心价值 & 不足 \\
\Block{3-1}{学术侧} &
\Block{3-1}{GPU 渲染、浏览器与微架构侧信道} &
\Block{3-1}{University of Texas at Austin、Carnegie Mellon University、University of Washington；Georgia Institute of Technology、University of Michigan、Ruhr University Bochum；University of California, Riverside 等浏览器/GPU 侧信道研究团队} &
GPU.zip 证明对软件透明的图形数据压缩可形成跨源像素泄露侧信道 \cite{gpuzip2024}。 &
利用 GPU 压缩、渲染延迟、像素布局差异构造跨源观测。 &
说明 AI Agent 的浏览器入口不只是应用层风险，也包含 GPU 渲染路径风险。 &
攻击效果依赖具体 GPU、驱动、浏览器和页面布局。 \\
& & &
Hot Pixels 展示频率、功耗、温度等混合信号可支持浏览器侧像素窃取和历史嗅探 \cite{hotpixels2023}。 &
通过 DVFS、温度、功耗传感器等软件可见信号推断图形内容。 &
推动浏览器、驱动和 GPU API 在隔离粒度、计时精度、压缩策略上重新评估。 &
缓解措施可能牺牲图形性能、压缩效率或 WebGPU 可用性。 \\
& & &
WebGPU-SPY 将 GPU cache 攻击引入 WebGPU 沙箱，显示浏览器 GPU 计算栈也可形成指纹识别面 \cite{webgpuispy2024}。 &
借助 WebGPU 高并行计时与 cache 争用建立浏览器内侧信道。 &
为端侧多模态 Agent、浏览器 Agent 的底层攻击面建模提供证据。 &
与大模型推理容器、企业浏览器环境之间还需要更多真实部署验证。 \\
\Block{3-1}{学术侧} &
\Block{3-1}{GPU 显存、本地内存与多租户 AI 算力隔离} &
\Block{3-1}{Trail of Bits、University of California, Santa Cruz；University of California, Riverside、Pacific Northwest National Laboratory 等 GPU 内存与多 GPU 侧信道研究团队} &
LeftoverLocals 证明部分 GPU local memory 清理不足可泄露跨应用、跨用户或跨容器数据，研究与协调披露涉及 AMD、Apple、Qualcomm 等厂商 \cite{leftoverlocals2024,certleftoverlocals2024}。 &
读取 GPU kernel 结束后残留的 local memory，恢复上一个任务的中间数据。 &
将 AI/ML 负载中的 prompt、响应、中间张量和模型片段纳入 GPU 隔离威胁模型。 &
攻击复现往往依赖特定设备、API、驱动版本和调度条件。 \\
& & &
Spy in the GPU-box 证明 NVIDIA DGX 多 GPU 系统中的远端 GPU L2 cache 争用可形成隐蔽信道和侧信道 \cite{spygpubox2022}。 &
逆向多 GPU cache/互连层次，通过 contention、prime-probe、性能计数器等方式观测远端负载。 &
说明共享 GPU 不是普通共享 CPU 资源，Agent 运行时需要单独建模显存、kernel、驱动和调度器。 &
内存清零、禁用并发、限制计数器会带来吞吐和利用率损失。 \\
& & &
NVBleed/Beyond the Bridge 等后续工作进一步把 NVLink/PCIe 互连纳入云端多 GPU 泄露面分析 \cite{nvbleed2025,beyondbridge2024}。 &
在多租户或虚拟化场景中测试互连拥塞、缓存命中和应用指纹泄露。 &
为云 GPU 分配、容器隔离、MIG/vGPU 策略和清零机制提供实证依据。 &
云厂商的真实多租户策略和硬件细节不完全公开，第三方验证仍有限。 \\
\Block{3-1}{学术侧} &
\Block{3-1}{硬件木马、供应链可信与安全 EDA} &
\Block{3-1}{University of Connecticut、Rice University；University of Florida；DARPA AISS 及安全 EDA/IP 生态相关团队} &
硬件木马分类与检测研究将威胁从制造环节扩展到 IP 复用、设计外包、测试不足和供应链可信 \cite{tehranipoor2010trojan}。 &
建立硬件木马 taxonomy、触发机制、载荷类型和检测方法。 &
回答 AI Agent 底座的更深层问题：芯片、IP、固件和供应链是否可信。 &
与具体 AI 加速器和大模型负载的结合仍不充分。 \\
& & &
信息流安全验证等方法尝试在 RTL/IP 层检测恶意逻辑或违反安全策略的设计 \cite{nahiyan2018ifs}。 &
在设计阶段使用信息流、形式化验证、测试向量和侧信道分析识别异常。 &
将安全从部署后补丁前移到芯片设计与验证阶段。 &
形式化验证和检测方法面对大规模 SoC 时成本高、覆盖有限。 \\
& & &
DARPA AISS 将侧信道、逆向工程、供应链攻击和恶意硬件纳入自动化安全芯片设计流程 \cite{darpaAISS}。 &
将安全分区、IP 完整性、设计空间评估和安全成本权衡融入 SoC 自动化流程。 &
为 AI 加速器、NPU、DPU 和定制 SoC 的安全认证提供基础方法。 &
供应链可信依赖厂商、代工、EDA、IP 提供商协同，落地复杂。 \\
\Block{3-1}{产业侧} &
\Block{3-1}{云端机密 AI 加速与 CPU-GPU TEE 协同} &
\Block{3-1}{NVIDIA Confidential Computing and Attestation；AMD SEV-SNP；Intel Trust Domain Extensions；云端机密计算部署平台} &
NVIDIA H100 将 GPU 纳入机密计算边界，提供硬件 RoT、secure/measured boot、SPDM、设备证明、CC-On 模式和加密 CPU-GPU 传输 \cite{nvidiah100cc2023}。 &
在 GPU 内建立硬件 RoT、固件签名、测量启动和设备证书链。 &
保护 AI 训练/推理中的数据、模型权重和企业上下文在“使用中”的机密性。 &
信任链跨 CPU、GPU、固件、驱动、证明服务和云编排，TCB 复杂。 \\
& & &
AMD SEV-SNP 和 Intel TDX 为机密 VM 提供内存加密、完整性保护和远程证明，可与 GPU 机密计算形成协同 \cite{amdsev2026,inteltdx2026}。 &
用 CVM 保护主机侧内存、完整性与远程证明，并与 GPU TEE 的设备证明和受保护数据通道形成协同。 &
支撑金融、医疗、政务和企业 Agent 在公有云/混合云中部署敏感 AI 负载。 &
PCIe 加密、数据搬运和证明流程可能带来性能与运维成本。 \\
& & &
NVIDIA Secure AI/NRAS 等产品化能力把 GPU 证明、密钥发布、镜像治理和企业 AI 工作流连接起来 \cite{nvidiasecureai2025}。 &
通过 NVIDIA 远程证明和策略校验能力验证 GPU、驱动、固件、镜像和运行模式后再释放敏感数据或模型。 &
把硬件安全能力转化为云服务和企业 AI 平台的可售卖能力。 &
关键实现仍高度依赖厂商闭源固件、证书服务和云平台策略。 \\
\Block{3-1}{产业侧} &
\Block{3-1}{GPU 漏洞响应、驱动修复与硬件安全治理} &
\Block{3-1}{Trail of Bits、CERT/CC、AMD Product Security、Apple、Google/Qualcomm 等相关厂商安全响应团队} &
CERT/CC VU\#446598 围绕 LeftoverLocals 对 GPU kernel/local memory 泄漏进行协调披露，涉及 AMD、Apple、Qualcomm 等多家 GPU/平台厂商 \cite{certleftoverlocals2024}。 &
通过驱动、固件、运行模式限制并发 GPU 进程并清理寄存器/local memory。 &
把学术/安全公司发现快速转化为厂商修复路线。 &
修复链条长，涉及芯片商、设备商、OS、驱动、云平台和用户升级。 \\
& & &
AMD-SB-6010 将该问题归入 GPU Memory Leaks，并给出 CVE-2023-4969、驱动/固件修复和可选新运行模式 \cite{amdleftoverlocals2024}。 &
以 CVE、供应商公告、标准组织和协调披露机制推动跨生态修复。 &
强化 AI 负载中 GPU local memory、寄存器、kernel 边界的安全基线。 &
可选隔离模式可能降低并发能力和 GPU 利用率。 \\
& & &
AMD、Apple 等厂商围绕受影响设备、驱动/固件版本和运行模式发布修复或缓解，Google、Qualcomm 等相关状态需以各自公告或 CERT vendor status 为准 \cite{certleftoverlocals2024}。 &
对 OpenCL、Metal、Vulkan、ROCm、系统驱动和设备固件进行版本化治理。 &
促使 GPU 厂商把多租户 AI、容器和 LLM 场景纳入安全公告。 &
老旧设备、移动 SoC 和闭源驱动的长期修复可达性不稳定。 \\
\Block{3-1}{产业侧} &
\Block{3-1}{端云协同私有 AI 与可验证云推理节点} &
\Block{3-1}{Apple Private Cloud Compute；Apple Silicon、Secure Enclave、Apple Intelligence 生态} &
Apple PCC 将 Apple Silicon、Secure Enclave、Secure Boot、代码签名和受限 OS 引入云端 AI 推理节点 \cite{applepcc2024}。 &
用端侧设备先验证 PCC 节点证书、系统测量值和透明日志，再加密发送请求。 &
将端侧隐私模型延伸到云端大模型推理，形成端云协同可信 AI 底座。 &
方案强依赖 Apple 封闭硬件、OS 和服务生态，跨平台可复用性有限。 \\
& & &
PCC 明确提出无远程 shell、无通用日志、无特权运行时访问和请求处理后不保留用户数据 \cite{applepcc2024}。 &
通过 Secure Boot、代码签名、trust cache、SEP 密钥保护和内存安全语言缩小攻击面。 &
从“云服务承诺不看数据”转向“硬件、系统镜像和证明链共同约束服务提供者”。 &
公开透明度虽高于传统云 AI，但关键硬件与部分实现仍由厂商控制。 \\
& & &
PCC 承诺通过透明日志、公开生产软件镜像和设备侧证明，让用户设备只向可验证节点发送请求 \cite{applepcc2024}。 &
通过无特权接口、受限观测、请求后清除和 target diffusion 降低运营方与定向攻击风险。 &
对移动端 Agent、个人智能助理和隐私敏感推理具有示范意义。 &
主要解决云端推理隐私，不直接处理提示注入、工具越权和长期日志治理。 \\
\end{NiceTabularX}
\caption{北美 AI 硬件与芯片安全研究现状：以方向为主轴、主体为支撑的结构化总结}
\label{tab:north-america-ai-hardware-security-nicematrix}
\end{table*}
